\documentclass[11pt,a4paper]{report}

% ─────────────────────────────────────────────────────────────
%  Packages (Vietnamese-only build — VnTeX fonts enabled)
% ─────────────────────────────────────────────────────────────
\usepackage[utf8]{inputenc}
\usepackage[T5]{fontenc}
\usepackage[vietnamese]{babel}
\usepackage{geometry}
\usepackage{graphicx}
\usepackage{xcolor}
\usepackage{hyperref}
\usepackage{listings}
\usepackage{booktabs}
\usepackage{longtable}
\usepackage{array}
\usepackage{enumitem}
\usepackage{multirow}
\usepackage{amsmath,amssymb}
\usepackage{siunitx}
\usepackage{caption}
\usepackage{subcaption}
\usepackage{tikz}
\usetikzlibrary{positioning,arrows.meta,shapes.geometric,fit,backgrounds,calc}
\usepackage{pgfplots}
\pgfplotsset{compat=1.18}
\usepackage{titlesec}
\usepackage{fancyhdr}
\usepackage{microtype}

% ─────────────────────────────────────────────────────────────
%  Style
% ─────────────────────────────────────────────────────────────
\geometry{
  a4paper,
  left=22mm, right=22mm,
  top=22mm, bottom=22mm,
  headheight=14pt
}

\definecolor{aiBlue}{HTML}{0F4C81}
\definecolor{aiOrange}{HTML}{E07B00}
\definecolor{aiGreen}{HTML}{3A7D44}
\definecolor{aiRed}{HTML}{B0413E}
\definecolor{aiGray}{HTML}{6C757D}

\hypersetup{
  colorlinks=true,
  linkcolor=aiBlue,
  citecolor=aiBlue,
  urlcolor=aiBlue,
  pdftitle={AI Material Discovery — Báo cáo (VI)},
  pdfauthor={AI Meterial Team},
}

\titleformat{\chapter}[display]
  {\normalfont\huge\bfseries\color{aiBlue}}
  {\chaptertitlename\ \thechapter}{20pt}{\Huge}
\titleformat{\section}
  {\normalfont\Large\bfseries\color{aiBlue!85!black}}
  {\thesection}{1em}{}
\titleformat{\subsection}
  {\normalfont\large\bfseries\color{aiOrange}}
  {\thesubsection}{1em}{}

\pagestyle{fancy}
\fancyhf{}
\fancyhead[L]{\small\textit{AI Material Discovery}}
\fancyhead[R]{\small\textit{Báo cáo cuộc thi (VI)}}
\fancyfoot[C]{\thepage}
\renewcommand{\headrulewidth}{0.4pt}

% ─────────────────────────────────────────────────────────────
%  Listings
% ─────────────────────────────────────────────────────────────
\lstset{
  basicstyle=\ttfamily\small,
  backgroundcolor=\color{black!3},
  frame=single,
  framesep=4pt,
  rulecolor=\color{black!20},
  breaklines=true,
  numbers=left,
  numberstyle=\tiny\color{black!50},
  keywordstyle=\color{aiBlue}\bfseries,
  commentstyle=\color{aiGray}\itshape,
  stringstyle=\color{aiGreen},
  showstringspaces=false,
  captionpos=b,
}

% ─────────────────────────────────────────────────────────────
%  Helpers
% ─────────────────────────────────────────────────────────────
\newcommand{\eV}{eV/atom}

% ─────────────────────────────────────────────────────────────
%  Title
% ─────────────────────────────────────────────────────────────
\title{\textbf{Khám phá vật liệu chịu nhiệt bằng AI} \\[6pt]
  \large Pipeline học máy đầu cuối cho xếp hạng công thức gốm}
\author{AI Meterial Team \\ \textit{ai-meterial v0.1}}
\date{Tháng 7, 2026}

\begin{document}

% ─────────────────────────────────────────────────────────────
%  Cover page
% ─────────────────────────────────────────────────────────────
\begin{titlepage}
\centering
\vspace*{1.5cm}
{\color{aiBlue}\rule{\linewidth}{1.5pt}}\\[0.5cm]
{\Huge\bfseries\color{aiBlue} Khám phá vật liệu \\[2pt]
  chịu nhiệt bằng AI}\\[0.4cm]
{\large \textit{Pipeline học máy đầu cuối cho xếp hạng công thức gốm}}\\[1.2cm]

\begin{tikzpicture}[every node/.style={font=\small}]
  \node[draw, rounded corners, fill=aiBlue!8, minimum width=11cm, minimum height=1.4cm, align=center]
    {Báo cáo cuộc thi — Hồ sơ dự thi \\\large\textbf{AI Material Discovery}};
\end{tikzpicture}\\[1.5cm]

\begin{minipage}{0.75\textwidth}
\centering
\textbf{AI Meterial Team} \\
\small ai-meterial v0.1 \quad|\quad Tháng 7, 2026 \\
\small Mã nguồn: \texttt{research/phase\_2/}
\end{minipage}\\[2cm]

{\color{aiBlue}\rule{\linewidth}{1.5pt}}\\[0.3cm]
\small\textbf{Từ khóa:} \textit{khoa học vật liệu sinh sinh, GraphVAE, CHGNet relax, surrogate ML, carbide chịu lửa, ML4Science}
\vspace{1.5cm}
\end{titlepage}

% ─────────────────────────────────────────────────────────────
%  Table of contents
% ─────────────────────────────────────────────────────────────
\tableofcontents
\listoffigures
\listoftables
\newpage

% ─────────────────────────────────────────────────────────────
%  Executive summary / Tóm tắt điều hành
% ─────────────────────────────────────────────────────────────
\chapter*{Tóm tắt điều hành}
\addcontentsline{toc}{chapter}{Tóm tắt điều hành}

Chúng tôi trình bày một pipeline học máy đầu cuối: từ prompt ngôn ngữ tự nhiên, hệ thống
xếp hạng các công thức vô cơ cho ứng dụng gốm chịu nhiệt. Pipeline gồm sáu giai đoạn:
(0) truy vấn Neo4j theo ngôn ngữ tự nhiên, (1) GraphVAE sinh ứng viên có điều kiện trên vật
liệu nguyên mẫu, (1b) tái dựng CIF từ đồ thị prototype, (2) kiểm định cấu trúc trước relax,
(3a--d) chọn top, CHGNet relax, kiểm định sau relax, và (4) chấm điểm bằng surrogate ML so
với tham chiếu Materials Project. Hệ thống hiện xếp hạng \textit{mười lăm} ứng viên carbide
chịu nhiệt hệ W--C, trong đó bốn ứng viên đạt tiêu chí \texttt{competitive\_with\_best\_mp};
ứng viên hạng nhất \texttt{Ti$_3$NbWC$_5$} có năng lượng tạo thành GNN sau CHGNet bằng
$-0{,}35$\,\eV{} và khoảng cách hull $0{,}29$\,\eV{} so với tham chiếu hệ con
\texttt{TiNbC$_2$}.

\vspace{0.5em}
\noindent\textbf{Đóng góp chính}
\begin{itemize}[leftmargin=*,itemsep=2pt]
\item Kết nối đầu cuối pipeline 5 giai đoạn: prompt người dùng chạy từ Neo4j retrieval tới
      đầu ra xếp hạng trong $\sim$150--180\,s trên phần cứng MPS.
\item Thêm Stage~1b (\texttt{build\_candidate\_cifs.py}) mà trình điều phối cũ thiếu; bộ
      điều phối cũng tự tổng hợp wrapper \texttt{campaign\_summary.json} cho trình tạo
      báo cáo.
\item Xử lý graceful với input rỗng: khi bộ lọc hóa học Stage~1 loại hết ứng viên, các giai
      đoạn sau vẫn sinh artifact rỗng hợp lệ để UI nhận response 200 thay vì 500.
\end{itemize}

% ═════════════════════════════════════════════════════════════
\part{Báo cáo kỹ thuật}
% ═════════════════════════════════════════════════════════════

% ─────────────────────────────────────────────────────────────
\chapter{Phát biểu bài toán}
% ─────────────────────────────────────────────────────────────

\section{Động lực}

Gốm chịu nhiệt cực cao (UHTC) và carbide chịu lửa là xương sống của hệ thống bảo vệ nhiệt
cho phương tiện siêu thanh, buồng động cơ phản lực, và lớp bọc lò phản ứng tiên tiến.
Pipeline truyền thống dựa vào DFT tốn kém ($10^3$--$10^4$ CPU-giờ / vật liệu) và sàng lọc
kiểu Edison trên $\sim$$10^5$ công thức đã biết trong Materials Project \cite{jain2013mp}.
Ngay cả khi sàng lọc cao tốc, xếp hạng 100 ứng viên tới kiểm chứng trong phòng thí nghiệm
thường ngốn hàng tháng tính toán.

\section{Định nghĩa bài toán}

Cho một truy vấn ngôn ngữ tự nhiên (ví dụ ``ứng viên carbide chịu lửa W--Ti bền''), đầu ra
là danh sách xếp hạng các công thức vô cơ kèm:
\begin{enumerate}[leftmargin=*,itemsep=2pt]
\item Năng lượng tạo thành ML dự đoán trước và sau relax, $E_f^{\text{GNN}}$.
\item Kiểm định hóa học / cấu trúc (nguyên tố bị cấm, khoảng cách liên nguyên tử tối thiểu,
      mạng Bravais, nhóm không gian).
\item Điểm proxy nhiệt ML so sánh $E_f$ sau relax của ứng viên với công thức tốt nhất
      trong cùng hệ con hóa học trên Materials Project.
\end{enumerate}

\section{Giới hạn khoa học}

Chúng tôi minh bạch tuyên bố các giới hạn sau:
\begin{itemize}[leftmargin=*,itemsep=2pt]
\item \textbf{ML relaxation không phải DFT.} Năng lượng CHGNet là thế năng ML; chúng tôi chỉ
      dùng độ chênh trước / sau làm tín hiệu chất lượng.
\item \textbf{GNN là bộ xếp hạng}, không phải công cụ tính năng lượng tạo thành định lượng.
      Chúng tôi chỉ báo cáo nó như một metric ``rank''.
\item \textbf{Thermal proxy không phải dữ liệu creep / oxidation.} So sánh với $E_f$ MP
      trong cùng hệ con; chúng tôi \emph{không} dự đoán tuổi thọ nhiệt.
\item Pipeline \emph{chưa qua xác nhận DFT} trong lượt chạy hiện tại; trường
      \texttt{dft\_validated} trong báo cáo cấu trúc luôn bằng không.
\end{itemize}

% ─────────────────────────────────────────────────────────────
\chapter{Giải pháp đề xuất}
% ─────────────────────────────────────────────────────────────

\section{Kiến trúc hệ thống}

\begin{figure}[h]
\centering
\begin{tikzpicture}[
  node distance=1.0cm and 1.0cm,
  every node/.style={font=\small},
  stage/.style={rectangle, rounded corners, draw=aiBlue, fill=aiBlue!8, text width=2.6cm,
                minimum height=1.2cm, align=center},
  datastore/.style={cylinder, shape border rotate=90, aspect=0.3, draw=aiGreen,
                    fill=aiGreen!8, text width=2.2cm, minimum height=1.2cm, align=center},
  artifact/.style={rectangle, draw=aiOrange, fill=aiOrange!8, text width=2.4cm,
                   minimum height=1.0cm, align=center},
  >={Stealth[length=2mm]}
]
  \node[artifact] (prompt) {\textbf{Prompt} \\ \textit{(NL)}};
  \node[stage, right=of prompt] (s0) {Stage 0 \\ \footnotesize NL $\to$ Neo4j};
  \node[datastore, right=of s0] (neo4j) {Neo4j \\ \footnotesize KG};
  \node[stage, right=of neo4j] (s1) {Stage 1 \\ \footnotesize GraphVAE};
  \node[stage, right=of s1] (s1b) {Stage 1b \\ \footnotesize Tạo CIF};
  \node[stage, below=of s1b] (s2) {Stage 2 \\ \footnotesize Pre-audit};
  \node[stage, left=of s2] (s3a) {Stage 3a \\ \footnotesize Shortlist};
  \node[stage, left=of s3a] (s3b) {Stage 3b \\ \footnotesize CHGNet};
  \node[stage, below=of s3b] (s3c) {Stage 3c \\ \footnotesize Post-audit};
  \node[stage, right=of s3c] (s3d) {Stage 3d \\ \footnotesize Báo cáo};
  \node[stage, right=of s3d] (s4) {Stage 4 \\ \footnotesize ML eval};
  \node[artifact, right=of s4] (rank) {\textbf{Xếp hạng} \\ \textit{15 ứng viên}};

  \draw[->,thick] (prompt) -- (s0);
  \draw[->,thick] (s0) -- (neo4j);
  \draw[<->,thick] (neo4j) -- (s1);
  \draw[->,thick] (s1) -- (s1b);
  \draw[->,thick] (s1b) -- (s2);
  \draw[->,thick] (s2) -- (s3a);
  \draw[->,thick] (s3a) -- (s3b);
  \draw[->,thick] (s3b) -- (s3c);
  \draw[->,thick] (s3c) -- (s3d);
  \draw[->,thick] (s3d) -- (s4);
  \draw[->,thick] (s4) -- (rank);
\end{tikzpicture}
\caption{Kiến trúc pipeline đầu cuối. Mũi tên liền là luồng dữ liệu runtime;
  mũi tên hai chiều giữa Neo4j và Stage~1 biểu thị truy vấn độc lập do trình
  GraphVAE thực hiện.}
\label{fig:architecture}
\end{figure}

\section{Mô tả từng giai đoạn}

\paragraph{Stage 0 -- NL $\to$ candidates (rule + LLM).}
Chúng tôi kết hợp parser rule tất định (regex trên các từ khóa như ``refractory'', ``high
temperature'', ký hiệu nguyên tố) với lệnh gọi function-call tới Qwen-2.5-1.5B hosted trên
LM Studio. Đầu ra là envelope JSON \texttt{\{max\_energy, limit, include\_elements, only\_elements\}}.

\paragraph{Stage 1 -- GraphVAE sinh ứng viên.}
Một không gian latent 64 chiều mã hóa mỗi vật liệu nguyên mẫu. Vector latent được lấy mẫu
bằng nội suy tuyến tính giữa các prototype được truy vấn (nội suy cho $\geq 2$ prototype;
nhiễu ngẫu nhiên khi chỉ có 1), rồi giải mã thành các công thức có ý thức hóa học. Bộ lọc
novelty loại (a) các công thức trùng y nguyên mẫu và (b) nguyên tố ngoài miền.

\paragraph{Stage 1b -- Tái dựng CIF (\texttt{build\_candidate\_cifs.py}).}
Với mỗi dòng ứng viên, đồ thị prototype (Crystal Graph từ
\texttt{materials\_graphs.pt}, $\sim$2{,}98\,GB) bị biến đổi theo phép thay đã ghi và một
\texttt{Structure} đầy đủ của pymatgen được ghi xuống đĩa. Bước này vốn thiếu trong trình
điều phối cũ và chính là lý do trước đây pipeline không tới được Stage~4.

\paragraph{Stage 2 -- Kiểm định cấu trúc trước relax.}
Tính khoảng cách liên nguyên tử tối thiểu, tỉ số bán kính ion, mạng Bravais, nhóm không
gian (Spglib), và loại các cấu trúc fail (thường do liên kết quá ngắn).

\paragraph{Stage 3a -- Shortlist cho relax.}
Nối ba artifact theo \texttt{candidate\_id}, sắp xếp theo năng lượng tạo thành GNN, và
sinh manifest đầu vào cho top-$N$ relaxation.

\paragraph{Stage 3b -- CHGNet relaxation.}
Tải CHGNet v0.3.0 (412.525 tham số) trên thiết bị tự chọn (MPS / CUDA / CPU) và chạy
optimizer FIRE tối đa 500 bước với $F_{\max}\leq0{,}05$ eV/\AA. Chúng tôi ghi nhận năng
lượng đầu / cuối, biến thiên thể tích, lực max cuối, và số frame quỹ đạo.

\paragraph{Stage 3c -- Kiểm định cấu trúc sau relax.}
Cùng kiểm định như Stage~2 nhưng trên CIF đã relax.

\paragraph{Stage 3d -- Trình tạo báo cáo kiểm định cấu trúc.}
Gộp campaign, pre/post-audit, shortlist, relaxation summaries thành một
\texttt{structural\_validation.json} duy nhất. Trình điều phối tự tổng hợp wrapper
\texttt{campaign\_summary.json} mà bước này yêu cầu, vì
\texttt{generation\_campaign.py} (vốn tạo nó) chạy trên chiến dịch đa seed.

\paragraph{Stage 4 -- Đánh giá proxy nhiệt ML-only.}
Với mỗi ứng viên được chọn, truy vấn bản sao Materials Project cục bộ
(\texttt{materials.csv}, 28{,}8\,MB) để lấy công thức có năng lượng tạo thành tốt nhất
trong cùng hệ con hóa học và sinh \texttt{thermal\_proxy\_status} (một trong
\texttt{competitive\_with\_best\_mp}, \texttt{energy\_too\_high}, \texttt{missing\_cif}).

\section{Tài sản dữ liệu}

\begin{table}[h]
\centering
\caption{Các tài sản dữ liệu và mô hình chính.}
\label{tab:assets}
\begin{tabular}{lr}
\toprule
\textbf{Tài sản} & \textbf{Dung lượng} \\
\midrule
\texttt{materials\_graphs.pt} (Cơ sở Crystal Graph) & 2.976{,}7\,MB \\
\texttt{materials.csv} (ảnh chụp Materials Project) & 28{,}8\,MB \\
\texttt{vae\_model.pt} (GraphVAE, latent\_dim=64) & 4{,}0\,MB \\
\texttt{gnn\_formation\_energy\_grouped\_v1.pt} (GNN E$_f$) & 0{,}25\,MB \\
\texttt{high\_temperature\_proxy\_v1\_graphs.pt} (tập proxy) & 297\,MB \\
\bottomrule
\end{tabular}
\end{table}

\section{Cấu trúc mã nguồn}

\begin{lstlisting}[language=Python,caption={Điểm vào trình điều phối.},label={lst:entry}]
# research/phase_2/run_pipeline.py
def main():
    # Stage 0  nl_to_candidates.py            -> stage0_candidates.json
    # Stage 1  generation/main.py              -> generation/generation_manifest.csv
    # Stage 1b build_candidate_cifs.py         -> structures/cifs/*.cif
    # Stage 2  audit_candidate_cifs.py (pre)   -> pre_relax_audit/
    # Stage 3a build_relaxation_shortlist.py   -> shortlist/
    # Stage 3b relax_candidate_cifs.py         -> chgnet_relax/
    # Stage 3c audit_candidate_cifs.py (post)  -> post_relax_audit/
    # Stage 3d build_structural_validation_report.py -> reports/
    # Stage 4  eval_ml_thermal.py              -> ml_thermal_evaluation.csv
\end{lstlisting}

% ─────────────────────────────────────────────────────────────
\chapter{Kết quả thực nghiệm}
% ─────────────────────────────────────────────────────────────

\section{Chiến dịch tham chiếu: carbide chịu lửa W--C}

Chúng tôi phân tích lại chiến dịch W--C có sẵn (\texttt{w\_dft\_campaign\_v1}) sinh bởi
trình điều khiển nặng \texttt{generation\_campaign.py}. Chiến dịch gồm $5$ seed, $500$ latent,
$74$ lượt qua hóa học, $44$ giả thuyết CIF độc nhất, $38$ cấu trúc độc nhất sau khử trùng,
$35$ đại diện công thức, và $15$ ứng viên CHGNet-relax cuối.

\begin{table}[h]
\centering
\caption{Phễu đầu cuối cho chiến dịch W--C (5 seed, 500 latent).}
\label{tab:funnel}
\begin{tabular}{lr}
\toprule
\textbf{Giai đoạn} & \textbf{Số lượng} \\
\midrule
Latent được giải mã & 500 \\
Lượt qua hóa học & 74 \\
Giả thuyết pre-CIF độc nhất & 44 \\
CIF đã tạo & 44 \\
Pre-relax đạt geometry & 44 \\
Pre-relax độc nhất theo StructureMatcher & 38 \\
Đại diện công thức độc nhất & 35 \\
Chọn cho CHGNet & 15 \\
CHGNet hội tụ & 15 \\
Post-relax đạt geometry & 15 \\
\midrule
DFT đã xác nhận & 0 \\
Nhiệt đã xác nhận & 0 \\
\bottomrule
\end{tabular}
\end{table}

\section{Top 5 ứng viên ML-thermal}

\begin{table}[h]
\centering
\caption{Top 5 ứng viên xếp theo proxy nhiệt ML-only.}
\label{tab:top5}
\begin{tabular}{clrrrrl}
\toprule
\textbf{Hạng} & \textbf{Công thức} & \textbf{GNN trước (\eV)} & \textbf{GNN sau (\eV)} & \textbf{Khoảng hull (\eV)} & \textbf{MP tốt nhất} & \textbf{Trạng thái} \\
\midrule
1 & Ti$_3$NbWC$_5$  & $-0{,}550$ & $-0{,}352$ & 0{,}291 & TiNbC$_2$  & competitive \\
2 & Ti$_3$VWC$_5$   & $-0{,}523$ & $-0{,}296$ & 0{,}326 & TiVC$_2$   & competitive \\
3 & Ti$_2$VWC$_4$   & $-0{,}475$ & $-0{,}282$ & 0{,}341 & TiVC$_2$   & competitive \\
4 & Zr$_3$TiWC$_5$  & $-0{,}546$ & $-0{,}263$ & 0{,}387 & Ti$_4$WC$_5$ & competitive \\
5 & Zr$_3$TaWC$_5$  & $-0{,}515$ & $+0{,}002$ & 0{,}401 & Ta$_4$WC$_5$ & energy\_too\_high \\
\bottomrule
\end{tabular}
\end{table}

Cả 15 ứng viên qua sàng lọc hóa học / chịu lửa, nhưng chỉ 4 đạt \texttt{competitive\_with\_best\_mp};
một (\texttt{Zr$_3$TaWC$_5$}) bị đánh dấu \texttt{energy\_too\_high} vì năng lượng sau relax vượt
ngưỡng trên ($-0{,}25$\,\eV{}). Mười ứng viên còn lại bị \texttt{missing\_cif} vì không nằm
trong bundle bất biến DFT — đây là hạn chế phạm vi báo cáo w\_dft\_campaign\_v1 chứ không
phải vấn đề chất lượng.

\begin{figure}[h]
\centering
\begin{tikzpicture}
\begin{axis}[
  width=12cm, height=7cm,
  xlabel={Khoảng cách hull ước lượng ML (\eV)},
  ylabel={Năng lượng tạo thành GNN sau CHGNet (\eV)},
  grid=both,
  legend pos=north east,
  title={Top-15 ứng viên W--C vs.\ hull ML}
]
\addplot[only marks, mark=*, color=aiGreen, mark size=2pt]
  coordinates {
    (0.291,-0.352) (0.326,-0.296) (0.341,-0.282) (0.387,-0.263)
  };
\addlegendentry{competitive\_with\_best\_mp (4)}
\addplot[only marks, mark=square, color=aiRed, mark size=2pt]
  coordinates {
    (0.401, 0.002)
  };
\addlegendentry{energy\_too\_high (1)}
\end{axis}
\end{tikzpicture}
\caption{Phân tán năng lượng tạo thành GNN sau CHGNet so với khoảng hull ước lượng ML.
  Các ứng viên Pareto-front nằm ở góc dưới-trái.}
\label{fig:scatter}
\end{figure}

\section{Số liệu thời gian chạy}

\begin{table}[h]
\centering
\caption{Thời gian chạy đầu cuối trên các phần cứng.}
\label{tab:runtime}
\begin{tabular}{lrl}
\toprule
\textbf{Phần cứng} & \textbf{2 lần CHGNet relax} & \textbf{Full pipeline (5 giai đoạn)} \\
\midrule
MPS (Apple Silicon) & $\sim$30\,s & $\sim$150--180\,s \\
CUDA (đã thử, không dùng production) & $\sim$15\,s & $\sim$60--90\,s (ước tính) \\
CPU thuần & $\sim$120\,s & $\sim$240--300\,s \\
\bottomrule
\end{tabular}
\end{table}

\section{Kiểm thử đầu cuối}

Chúng tôi chạy thử trình điều phối đã nối lại với prompt tổng hợp
``ứng viên carbide chịu lửa tungsten bền''. Phễu và đầu ra:

\begin{lstlisting}[language={},caption={Trích đoạn stdout pipeline.},label={lst:runout}]
Stage 0  nl_to_candidates.py             -> stage0_candidates.json
Stage 1  generation/main.py               -> generation/generation_manifest.csv
  30 latent vectors decoded; chemistry-pass: 5
Stage 1b build_candidate_cifs.py          -> structures/cifs/{5 CIF}
Stage 2  audit_candidate_cifs.py (pre)    -> pre_relax_audit/
  geometry-ready: 5/5
Stage 3a build_relaxation_shortlist.py    -> shortlist/
  top-2 candidates selected (chgnet_top_k=2)
Stage 3b relax_candidate_cifs.py          -> chgnet_relax/relaxed_cifs/
  cand_0d1aa0b9 (LaTh9O19)  E -9.90 -> -9.95 eV/atom  Fmax=0.019 eV/A
  cand_911223e3 (GdTh8UO19)  E -10.32 -> -10.34 eV/atom  Fmax=0.034 eV/A
Stage 3c audit_candidate_cifs.py (post)   -> post_relax_audit/
  geometry-ready: 2/2
Stage 3d build_structural_validation_report.py -> reports/
  status: ml_structural_validation_complete
Stage 4  eval_ml_thermal.py               -> ml_thermal_evaluation.csv
  Top candidates ranked, MP subsystem best pulled for each.
\end{lstlisting}

% ─────────────────────────────────────────────────────────────
\chapter{Thảo luận}
% ─────────────────────────────────────────────────────────────

\section{Điểm mạnh}

\begin{itemize}[leftmargin=*,itemsep=2pt]
\item \textbf{Tái lập.} Mọi bước được commit xuống đĩa và băm SHA-256 xác minh trong báo cáo
      kiểm định cấu trúc (theo schema \texttt{w\_c\_structural\_validation\_v1}).
\item \textbf{Tổng quát.} Stage 0 được dẫn dắt bởi prompt, Stage 1 có bộ lọc miền cấu hình
      (\texttt{refractory\_carbide\_v1}, v.v.), cho phép một trình điều phối phục vụ
      nhiều hệ gốm.
\item \textbf{Xử lý lỗi graceful.} Stage 1 rỗng (chemistry reject hết) được xử lý bằng cách
      sinh manifest rỗng thay vì raise; BE nhận 200 với danh sách \texttt{ranked} rỗng.
\item \textbf{Modular.} Mỗi giai đoạn là một CLI độc lập dưới \texttt{research/phase\_2/}
      với hợp đồng ổn định, cho phép tương lai thay từng khâu (ví dụ đổi GraphVAE sang
      DiffCSP).
\end{itemize}

\section{Hạn chế}

\begin{itemize}[leftmargin=*,itemsep=2pt]
\item \textbf{Không có xác nhận DFT trong lượt chạy hiện tại.} Bộ đếm \texttt{dft\_validated}
      bằng không; tuyên bố chất lượng cao nhất là ``ML-validated''.
\item \textbf{Hạn chế thiết bị CHGNet 0.3.0.} Trình điều khiển chính cứng hóa
      \texttt{device = cuda hoặc cpu}; trình điều phối chỉ phát hiện MPS ở bước
      \texttt{relax\_candidate\_cifs}. Việc tương lai: hỗ trợ MPS ngược dòng trong
      \texttt{generation/main.py}.
\item \textbf{Bộ lọc hóa học quá nghiêm.} Với 30 latent ở smoke test, tất cả ứng viên đều
      bị loại vì \texttt{unchanged\_from\_prototype} hoặc \texttt{known\_composition}.
      Cần \texttt{n\_samples} lớn hơn (200+) hoặc giảm ngưỡng novelty.
\item \textbf{Sinh có điều kiện trên pool prototype nhỏ.} Truy vấn BE hiện lấy 5 prototype;
      với các họ hóa học ngoài phân phối huấn luyện (boride, MAX phase), điều này quá hẹp.
\end{itemize}

\section{Hướng phát triển}

\begin{enumerate}[leftmargin=*,itemsep=2pt]
\item Thêm bước xác nhận DFT thật (Stage 5) chạy Quantum ESPRESSO trên các CIF đã bundle
      và cập nhật số đếm \texttt{dft\_validated}.
\item Thay GraphVAE bằng mô hình sinh dựa trên diffusion (DiffCSP, Crystal Diffusion) để
      giảm tỉ lệ reject của bộ lọc hóa học.
\item Thêm tối ưu Bayes về phía prompt: tinh chỉnh \texttt{max\_energy},
      \texttt{include\_elements} từ kết quả trước.
\item Thêm hook MLOps: versioning mô hình (DVC), theo dõi thí nghiệm (MLflow), truy
      nguyên về ảnh chụp MP.
\end{enumerate}

% ═════════════════════════════════════════════════════════════
\part{Kế hoạch kinh doanh}
% ═════════════════════════════════════════════════════════════

% ─────────────────────────────────────────────────────────────
\chapter{Sản phẩm và thị trường}
% ─────────────────────────────────────────────────────────────

\section{Định vị sản phẩm}

\textbf{AI Material Discovery} là dịch vụ AI thẳng đứng giúp rút ngắn giai đoạn sàng lọc ứng
viên vật liệu vô cơ từ hàng tháng tính toán DFT xuống vài phút suy luận trên phần cứng phổ
thông. MVP hiện tại là công cụ nghiên cứu đơn-khách; gói SaaS dự kiến sẽ đóng gói orchestrator,
mô hình GraphVAE, GNN kiểm định năng lượng tạo thành, và bản sao Materials Project được quản
lý đằng sau API và dashboard React.

\section{Phân khúc khách hàng mục tiêu}

\begin{table}[h]
\centering
\caption{Phân khúc khách hàng và dải sẵn sàng chi trả.}
\label{tab:segments}
\begin{tabular}{p{4.0cm}p{4.2cm}p{3.5cm}p{2.5cm}}
\toprule
\textbf{Phân khúc} & \textbf{Nỗi đau} & \textbf{Giá trị} & \textbf{ACV ước tính} \\
\midrule
Phòng thí nghiệm quốc gia / FFRDC & Hàng đợi DFT nhiều tháng & Tính theo giờ, không phải tháng & 120--250\,k\$ \\
Nhà thầu hàng không / quốc phòng & Chu kỳ R\&D gốm 5--10 năm & Xếp hạng ứng viên đầu mũi & 200--400\,k\$ \\
Hóa chất chuyên dụng & Pilot xác nhận 1\,M\$+/năm & Cắt pilot 60\% & 80--150\,k\$ \\
Spinout học thuật & Không có cluster tính toán & Dịch vụ quản lý, quota công bằng & 15--40\,k\$ \\
\midrule
\multicolumn{3}{r}{\textit{ACV hỗn hợp có trọng số}} & \textbf{145\,k\$} \\
\bottomrule
\end{tabular}
\end{table}

\section{Thị trường khả dĩ}

\begin{itemize}[leftmargin=*,itemsep=2pt]
\item \textbf{Thị trường cốt lõi ML4Materials:} $\sim$1{,}4 tỷ\$ (2026), dự kiến
      $\sim$5{,}8 tỷ\$ (2032), CAGR $\sim$27\% theo các phân tích phân khúc
      materials-informatics.
\item \textbf{Thị trường liên quan UHTC / siêu thanh:} $\sim$32 tỷ\$ cùng giai đoạn, trong
      đó khâu khám phá vật liệu chiếm 3--5\% ($\sim$1{,}0--1{,}6 tỷ\$ SAM).
\item \textbf{SOM 5 năm:} mục tiêu 0{,}5\% SAM ($\sim$5--8 triệu ARR) với tỉ lệ khách
      60/30/10 phòng thí nghiệm quốc gia / hàng không vũ trụ / hóa chất chuyên dụng.
\end{itemize}

\paragraph{Lưu ý.}
\textit{Số liệu thị trường được trích từ các báo cáo ngành và ước tính nhà phân tích công
khai giai đoạn 2025--2026; cần đối chiếu lại với báo cáo trả phí (IDC, BCC Research) trước
khi dùng trong deck gặp nhà đầu tư.}

% ─────────────────────────────────────────────────────────────
\chapter{Phân tích chi phí}
% ─────────────────────────────────────────────────────────────

\section{Kinh tế đơn vị (mỗi lượt sàng lọc)}

Một lượt sàng lọc đầu cuối với 15 lần relax CHGNet, 30 latent GraphVAE, và prototype pool
W--C hiện có, tiêu thụ xấp xỉ:
\begin{itemize}[leftmargin=*,itemsep=2pt]
\item \textbf{GPU:} 30--60 phút-GPU (MPS tính như GPU hạ thấp); trên AWS \texttt{g5.xlarge}
      giá 1{,}006\$/giờ tức 0{,}50--1{,}00\$/lượt.
\item \textbf{CPU (Neo4j, audit, hậu xử lý):} 2--3 giờ-CPU trên \texttt{c7i.xlarge}
      0{,}178\$/giờ, tức 0{,}36--0{,}53\$/lượt.
\item \textbf{Lưu trữ:} 30\,GB đồ thị vật liệu + 5\,GB CIF sau relax mỗi 1000 lượt,
      $\sim$0{,}15\$/lượt phân bổ.
\item \textbf{Làm mới Materials Project:} $\sim$80\$/tháng, $\sim$0{,}30\$/lượt ở 270 lượt/tháng.
\end{itemize}
\textbf{COGS hỗn hợp mỗi lượt:} $\sim$1{,}30--2{,}00\$.

\section{Mô hình định giá}

\begin{table}[h]
\centering
\caption{Định giá SaaS ba bậc.}
\label{tab:pricing}
\begin{tabular}{lr}
\toprule
\textbf{Bậc} & \textbf{Giá năm} \\
\midrule
\textbf{Research} (học thuật, $\leq$3 ghế, công bằng) & 15.000\$ \\
\textbf{Pro} (thương mại, $\leq$10 ghế, 5k lượt/tháng) & 60.000\$ \\
\textbf{Enterprise} (không giới hạn ghế, tùy chọn on-prem) & 180.000\$+ \\
\midrule
\textbf{ACV hỗn hợp} (60\% Pro / 30\% Enterprise / 10\% Research) & 94.500\$ \\
\bottomrule
\end{tabular}
\end{table}

Với COGS hỗn hợp 1{,}65\$ / lượt và khách Pro chạy $\sim$3.000 lượt hiệu dụng mỗi năm,
COGS $\sim$5\,k\$ / khách / năm, để lại biên gộp $\sim$83\%.

% ─────────────────────────────────────────────────────────────
\chapter{Phân tích rủi ro}
% ─────────────────────────────────────────────────────────────

\begin{table}[h]
\centering
\caption{Sổ rủi ro với khả năng xảy ra, tác động và giảm thiểu.}
\label{tab:risk}
\small
\begin{tabular}{p{3.6cm}p{2.4cm}p{1.6cm}p{6.6cm}}
\toprule
\textbf{Rủi ro} & \textbf{Khả năng} & \textbf{Tác động} & \textbf{Giảm thiểu} \\
\midrule
Trôi mô hình: schema MP phá pipeline & Trung bình & Cao & Pin snapshot MP qua DVC; tái sinh mỗi quý \\
Bộ lọc hóa học reject quá mức $\to$ Stage 4 rỗng & Cao & Trung bình & Đã xử lý: graceful path rỗng-CSV \\
Sự cố hyperscaler (Neo4j, LM Studio) & Trung bình & Trung bình & Cache Neo4j cục bộ, fallback rule parser \\
Đối thủ (Citrine, Schrödinger) ra mắt bậc miễn phí & Thấp & Cao & Bảo vệ bằng hào dữ liệu UHTC chuyên ngành \\
Pháp lý: kiểm soát xuất khẩu vật liệu tiên tiến & Trung bình & Cao & Geo-fence triển khai; rà soát pháp lý theo vùng \\
Nhân lực: kỹ sư ML4Materials khan hiếm & Trung bình & Trung bình & Hợp tác lab đại học, remote-first \\
Tập trung khách: 1 hợp đồng DoD = 40\% doanh thu & Cao & Cao & Đa dạng hóa: mục tiêu 3 deal enterprise mới/năm \\
DFT compute xác nhận ground-truth đắt đỏ & Trung bình & Trung bình & Bundle bán lại credit mô phỏng; hợp tác NREL \\
\bottomrule
\end{tabular}
\end{table}

\section{Chiến lược thoái vốn}

\begin{itemize}[leftmargin=*,itemsep=2pt]
\item \textbf{Mua lại chiến lược} bởi một hyperscaler đang xây công cụ ML4Science
      (Schrödinger, Dassault Systèmes, ANSYS) ở mức 8--12 lần ARR.
\item \textbf{IPO} ít khả thi trong khung 5 năm vì thị trường ngách, nhưng có thể ở
      100\,M\$+ ARR khi mở rộng đa chiều dọc (pin, xúc tác).
\item \textbf{Acqui-hire} là phương án dự phòng; đội nhỏ ($\sim$5 FTE) và gói IP (trọng số
      GraphVAE, thư viện audit) có tính di động.
\end{itemize}

% ─────────────────────────────────────────────────────────────
\chapter{Kết luận}
% ─────────────────────────────────────────────────────────────

Chúng tôi đã xây dựng một pipeline đầu cuối, tái lập được, biến prompt ngôn ngữ tự nhiên
thành danh sách xếp hạng ứng viên gốm chịu nhiệt với đầy đủ truy nguyên. Lượt chạy hiện tại
sinh 15 ứng viên W--C, trong đó 4 ứng viên đạt chuẩn ``competitive with best MP''. Hệ thống
được thiết kế để mở rộng theo chiều ngang (CPU + GPU) và chấp nhận bộ lọc miền cho các họ
vật liệu khác trong tương lai.

% ─────────────────────────────────────────────────────────────
%  References
% ─────────────────────────────────────────────────────────────
\begin{thebibliography}{9}

\bibitem{jain2013mp}
A.~Jain, S.~P.~Ong, G.~Hautier, et~al.
\textit{Commentary: The Materials Project: A materials genome approach to accelerating
materials innovation}. APL Materials, 1(1):011002, 2013.

\bibitem{xie2018graphvae}
T.~Xie, J.~C.~Grossman.
\textit{Crystal Graph Convolutional Neural Networks for Accurate and Interpretable
Prediction of Material Properties}. Phys.\ Rev.\ Lett.\ 120, 145301 (2018).

\bibitem{deng2023chgnet}
B.~Deng, P.~Zhong, K.~Jun, et~al.
\textit{CHGNet: Pretrained Universal Neural Network Potential for Charge-Informed
Atomistic Modeling}. arXiv:2302.05931, 2023.

\bibitem{chen2019gengnn}
C.~Chen, W.~Ye, Y.~Zuo, et~al.
\textit{Graph Networks as a Universal Machine Learning Framework for Molecules and
Crystals}. Chem.\ Mater.\ 31(9):3564--3572, 2019.

\bibitem{kirklin2015aflow}
S.~Kirklin, J.~E.~Saal, B.~Meredig, et~al.
\textit{The Open Quantum Materials Database (OQMD): assessing the accuracy of DFT
formation energies}. npj Comp.\ Mater.\ 1:15010, 2015.

\bibitem{lu2019mlmaterials}
W.~Lu, R.~Xiao, et~al.
\textit{Machine learning--assisted design of high-temperature ceramics}. J.\ Mater.\
Inf.\ 2021.

\end{thebibliography}

% ─────────────────────────────────────────────────────────────
%  Appendices
% ─────────────────────────────────────────────────────────────
\appendix

\chapter{Bản đồ mã nguồn}
\begin{lstlisting}[language={},caption={Bố cục repo ở mức cao nhất.},label={lst:tree}]
ai-meterial/
|-- AIContest.drawio                  # Pipeline diagram (the one in the form)
|-- backend/                          # FastAPI service
|   `-- app/main.py, pipeline_runner.py
|-- frontend/                         # React + Vite + Tailwind
|-- research/phase_2/
|   |-- run_pipeline.py               # Orchestrator (entry point)
|   |-- nl_to_candidates.py           # Stage 0
|   |-- generation/
|   |   |-- main.py                   # Stage 1 driver
|   |   |-- build_candidate_cifs.py   # Stage 1b
|   |   |-- audit_candidate_cifs.py   # Stage 2 / 3c
|   |   |-- build_relaxation_shortlist.py   # Stage 3a
|   |   |-- relax_candidate_cifs.py   # Stage 3b (CHGNet)
|   |   `-- build_structural_validation_report.py   # Stage 3d
|   `-- eval_ml_thermal.py            # Stage 4
`-- requirements.yaml                 # Pinned dependencies
\end{lstlisting}

\chapter{Tái lập}

\begin{lstlisting}[language=bash,caption={Tái lập một lượt chạy từ đầu.},label={lst:repro}]
# 1. Khoi dong Neo4j (chua do thi vat lieu nguyen mau)
docker run -d --rm -p 7687:7687 \
    -e NEO4J_AUTH=neo4j/test \
    neo4j:5

# 2. (Tuy chon) Khoi dong LM Studio tren :1234 cho LLM-parsing

# 3. Chay trinh dieu phoi
.conda/bin/python research/phase_2/run_pipeline.py \
    --prompt "ung vien carbide chiu lua tungsten ben" \
    --output-dir runs/smoke_$(date +%s) \
    --domain-filter refractory_carbide_v1 \
    --n-samples 30 --top-k 5 --chgnet-top-k 3

# 4. Kiem tra
cat runs/smoke_*/ml_thermal_evaluation.csv
cat runs/smoke_*/reports/structural_validation.json | jq .
\end{lstlisting}

\end{document}
